
Trustworthiness evaluation frameworks for large language models (LLMs) implicitly assume that safety, fairness, truthfulness, adversarial robustness, privacy, and machine ethics are orthogonal axes --- each measuring a distinct property that demands independent assessment. This assumption is operationalized in every major benchmark: TrustLLM \citep{Sun2024TrustLLM} evaluates 16 models across 6 dimensions, HELM \citep{Liang2022} tracks 7 metrics across 30 models, and MultiTrust \citep{Zhang2024MultiTrust} presents 5 dimensions as a radar chart. Yet no study has asked the empirical question underlying all of this evaluation: \textit{are these dimensions actually independent?}

We show they are not, at least in the TrustLLM 16-model evaluation setting. After controlling for model scale and RLHF alignment status, 8 of 15 trustworthiness dimension pairs are statistically correlated at Bonferroni-corrected significance, with partial Spearman coefficients reaching $\rho$ = 0.971 for the safety--privacy pair. The correlation structure is not noise --- it reflects a geometry driven by RLHF preference training that simultaneously shapes 5 of 6 trustworthiness dimensions, while adversarial robustness follows an independent trajectory.

This finding has two immediate practical implications. First, \textbf{evaluation compression}: a 3-dimension minimum evaluation set {truthfulness, fairness, privacy}, derived from the minimum spanning tree (MST) of the partial correlation matrix, spans the correlated cluster with mean bootstrap edge frequency 0.917 --- a 50\% reduction from the standard 6-dimension protocol. Second, \textbf{differential attention}: adversarial robustness cannot be inferred from the 5-dimension RLHF-shaped cluster and requires separate evaluation, as it follows a trajectory governed primarily by model scale rather than alignment.

The core mechanistic account is as follows: RLHF human preference annotation creates a shared optimization signal that annotators apply when rating outputs for harmlessness, fairness, and appropriate information handling. This shared signal simultaneously shapes safety, ethics, fairness, privacy, and truthfulness through the same preference-ranking process. Adversarial robustness --- which is not directly probed in standard preference annotation --- escapes this shared signal.

A notable finding that contradicts the original hypothesis is privacy's behavior. Privacy was predicted to be RLHF-insensitive (similar to robustness), but instead shows the strongest co-movement with safety of any pair ($\rho$ = 0.971). Two complementary explanations are plausible: RLHF reward models explicitly penalize privacy-violating outputs alongside harmful outputs, and TrustLLM's privacy dimension captures output-level refusal to produce PII, which mechanically overlaps with safety refusal behaviors.

We make the following contributions:

\begin{enumerate}
\item \textbf{The first, to our knowledge, systematic partial Spearman correlation analysis of LLM trustworthiness} --- We compute the full $6\times 6$ partial Spearman correlation matrix for TrustLLM's 16-model $\times$ 6-dimension evaluation data, controlling for model scale and RLHF status, revealing a strong 2-cluster structure (silhouette = 0.637, robust across all linkage methods) not previously reported.
\end{enumerate}

\begin{enumerate}
\item \textbf{Robustness isolation principle} --- Adversarial robustness is empirically identified as the sole RLHF-insensitive trustworthiness dimension. Within-family LLaMA-2 comparisons show negative robustness deltas after RLHF fine-tuning ($\Delta = -0.080, -0.085, -0.094$ at 7B/13B/70B), consistent with a directional safety-robustness tension, though this does not reach Bonferroni significance at n = 16.
\end{enumerate}

\begin{enumerate}
\item \textbf{RLHF joint optimization of safety and machine ethics} --- A within-family natural experiment on LLaMA-2 (7B, 13B, 70B) shows that RLHF fine-tuning simultaneously increases safety (mean $\Delta$ = +0.639) and machine ethics (mean $\Delta$ = +0.424) at every scale, providing directional mechanistic evidence for the 5-dimension RLHF-shaped cluster.
\end{enumerate}

\begin{enumerate}
\item \textbf{MST-derived minimum evaluation set} --- The minimum spanning tree identifies {truthfulness, fairness, privacy} as a sufficient 3-dimension evaluation set with mean bootstrap edge frequency 0.917 \citep{Tumminello2007MST}, enabling principled evaluation compression.
\end{enumerate}

The remainder of this paper is organized as follows: Section 2 surveys related work. Section 3 describes methodology. Section 4 details the experimental setup. Section 5 presents results. Section 6 discusses implications and limitations. Section 7 concludes.

